\chapter*{Conclusion}
\addcontentsline{toc}{chapter}{Conclusion}
\markboth{Conclusion}{Conclusion}

This thesis presented the design, development, and thorough performance evaluation of the \texttt{MLWE-Fast} library, a high-performance software implementation of the CRYSTALS-Kyber key encapsulation mechanism. Written in standard-compliant C11, the project was structured to separate key cryptographic layers, enabling targeted performance enhancements while preserving the mathematical boundaries required for security. By prioritizing portable C optimizations, such as division-free reduction algorithms and stack-allocated hash structures, the implementation demonstrates that significant speedups can be achieved over the standard reference codebase without sacrificing portability. This work highlights the critical intersection between abstract cryptographic specifications and physical microarchitectural execution, showing that software engineering decisions at the lowest level of arithmetic directly govern the efficiency of post-quantum communication.

The transition to post-quantum cryptography is necessitated by the rapid advancement of quantum computing hardware, which threatens the mathematical assumptions underlying traditional public-key schemes. With algorithms such as Shor's algorithm capable of solving integer factorization and discrete logarithms in polynomial time, widely deployed systems like RSA and elliptic curve cryptography will become obsolete. In response to this threat, standardizing bodies have focused on lattice-based cryptography, which is believed to be hard against both classical and quantum adversaries. By implementing the Kyber-512 parameter profile of Module-LWE, this thesis addresses the immediate need for efficient, software-based quantum-resistant implementations that can be integrated into existing network protocols without introducing prohibitive latency.

The software architecture of \texttt{MLWE-Fast} was designed with a three-layer structure that mirrors the algebraic design of the Module-LWE problem. The key encapsulation mechanism layer handles the public API and coordinates the Fujisaki-Okamoto transform, ensuring active security against chosen-ciphertext attacks by verifying ciphertexts through an explicit re-encryption loop. Beneath this wrapper, the public-key encryption layer coordinates the core matrix-vector operations, noise additions, and lossy compression steps. The arithmetic layer at the foundation executes the polynomial operations, Number Theoretic Transforms, and centered binomial distribution sampling. Correctness was verified at every step using a differential testing harness that executed key generation, encapsulation, and decapsulation operations in parallel with the official CRYSTALS-Kyber reference implementation, ensuring byte-for-byte compatibility.

The performance evaluation conducted on the AMD Zen 4 architecture yielded key insights into the behavior of cryptographic software under modern compilers and hardware. The most significant result was the massive speedup achieved by the portable Base C variant of the library, which ran $1.60\times$ faster for key pair generation, $2.01\times$ faster for encapsulation, and $2.07\times$ faster for decapsulation than the official reference implementation. Because this version contains no manual assembly or vector intrinsics, these performance gains are driven entirely by clean algorithmic design. Replacing division-heavy modular reduction steps with inlineable, division-free Montgomery and Barrett reductions eliminated high-latency division instructions from the critical path, allowing the CPU pipeline to execute arithmetic operations without instruction stalls.

In addition to reduction optimizations, the memory layout of the library was structured to maximize cache friendliness and align with the hardware boundaries of the target platform. By enforcing a strict 32-byte alignment on polynomial coefficient arrays, the data structures fit cleanly within the L1 and L2 cache lines of the Zen 4 CPU, avoiding split-line cache access penalties during sequential loops. Furthermore, by keeping intermediate variables stack-allocated and avoiding dynamic memory allocation within the main execution paths, the implementation minimized kernel-level overhead and preserved local CPU cache state. These engineering choices allowed the compiler to perform aggressive loop unrolling and register allocation, leading to highly optimized machine code under standard compiler flags.

While the Base C variant demonstrated substantial improvements over the reference code, the empirical evaluation of the hand-crafted AVX2 and AVX-512 variants revealed a surprising trend. Activating these manual SIMD code paths yielded only a marginal speedup of $2\%\text{--}3\%$ over the Base C implementation, and in some cases resulted in negligible performance decreases. This behavior highlights the limits of targeting manual optimization at isolated operations. Because the hand-written vector intrinsics were restricted to element-wise polynomial addition and subtraction, they only accelerated a tiny fraction of the total execution time. The bulk of the computation was dominated by symmetric hash evaluations and Number Theoretic Transforms, which remained scalar in all configurations, leaving the main bottlenecks untouched.

This performance profile is a classic validation of Amdahl's Law, which states that the speedup of a system from optimizing a single component is bounded by the fraction of time that the component is actually utilized. In the CRYSTALS-Kyber protocol, polynomial addition and subtraction constitute less than $5\%$ of the total execution cycles. Consequently, even an infinite acceleration of these addition loops can only yield a maximum theoretical speedup of approximately $5\%$. The primary computational cost lies in the Keccak-f[1600] permutation, which absorbs over $50\%$ of the runtime during matrix generation and Centered Binomial Distribution sampling, and the Cooley-Tukey butterfly kernels of the NTT, which consume $25\%\text{--}30\%$ of the cycles. Because these larger components remained scalar, the vectorization of additions had a negligible effect on the overall system performance.

The microarchitectural implementation of AVX-512 on the AMD Zen 4 processor further explains why the 512-bit vector registers did not double the performance of the 256-bit AVX2 variant. Unlike Intel architectures that utilize dedicated 512-bit wide execution units, Zen 4 schedules 512-bit vector instructions across its existing 256-bit pipelines, split over successive cycles. While this double-pumped design is highly efficient and prevents the severe frequency downclocking associated with dedicated 512-bit hardware, it means that the raw arithmetic throughput of a 512-bit instruction is comparable to executing two 256-bit instructions. As a result, the AVX-512 addition loops run at similar speeds to the AVX2 variants, and the compiler-optimized Base C loop (which is auto-vectorized to AVX2/AVX-512 by the compiler) achieves near parity.

These findings define a clear direction for future research and optimization within the library. To achieve higher levels of performance, the vectorization efforts must be expanded to target the major bottlenecks identified in the profiling analysis. The most promising path is the vectorization of the Number Theoretic Transform and inverse NTT kernels. By using 256-bit or 512-bit vector registers to execute Cooley-Tukey and Gentleman-Sande butterflies in a double-butterfly configuration, multiple coefficient updates can be performed in a single clock cycle. This can be achieved by loading precomputed twiddle factors into vectors and using register-shuffling instructions to permute coefficients within the vector registers, eliminating the need to write intermediate values back to cache.

Equally critical is the parallelization of the symmetric hashing layer, which represents the largest computational block in the protocol. Since matrix generation requires independent columns of the public matrix to be generated from a seed, a vectorized implementation can execute multiple independent Keccak-f[1600] permutations in parallel using wide SIMD registers. By running a four-way or eight-way parallel SHAKE128 generator, the library can generate multiple polynomial matrix entries simultaneously, fully utilizing the vector pipelines and lifting the performance ceiling. Combining parallel symmetric hashing with a vectorized NTT kernel will allow the library to achieve optimal performance on modern processors.

The development of the \texttt{MLWE-Fast} library highlights the importance of matching cryptographic algorithms with hardware execution characteristics. The project demonstrates that portable C code, when designed with cache alignment and division-free reduction algorithms, can achieve excellent performance that rivals hand-crafted assembly. However, it also shows that vectorization must be applied holistically to the entire execution path, rather than to minor arithmetic loops, to yield significant system-level benefits. By analyzing these microarchitectural limits and establishing the foundational pathways for NTT and Keccak vectorization, this thesis provides a solid engineering basis for the design of next-generation, high-performance post-quantum cryptographic software.